% Problemas del capítulo 10. Generado por tools/problemas/gen.py; no editar a mano.
\section{Problemas}\label{sec:problemas-10}

Graduados: \({\star}\) directo, \({\star}{\star}\) para pensar, \({\star}{\star}{\star}\) desafío. Las respuestas están en el apéndice~\ref{ap:respuestas}; conviene intentarlos antes de mirarlas.

\begin{problema}[\({\star}\)]\label{pb:10-1}
Calculá el jacobiano de \((r,\theta)\mapsto(r\cos\theta,\ r\sin\theta)\) en \(r=10\), \(\theta=30^\circ\), y su determinante. ¿Qué dice ese número sobre las áreas?
\end{problema}

\begin{problema}[\({\star}{\star}\)]\label{pb:10-2}
Desde una estación, un punto está en \(\Delta E=-120\)~m, \(\Delta N=50\)~m. Calculá la distancia y el acimut. Si el punto se corre \(1\)~cm al este, ¿cuánto cambian la distancia y el acimut? Usá el jacobiano de \(\mathrm{Pol}\).
\end{problema}

\begin{problema}[\({\star}{\star}\)]\label{pb:10-3}
Aproximá \(\sqrt{x^2+y^2}\) en \((3{,}02;\ 3{,}97)\) con su aproximación lineal en \((3,4)\) y compará con el valor exacto.
\end{problema}

\begin{problema}[\({\star}{\star}\)]\label{pb:10-4}
Un terreno tiene cota \(z=100-0{,}002\,x^2-0{,}001\,y^2\). En el punto \((100,50)\), ¿hacia dónde corre el agua y con qué pendiente?
\end{problema}

\begin{problema}[\({\star}{\star}{\star}\)]\label{pb:10-5}
Para \(h(x,y)=x^2+3xy+y^2\), calculá la hessiana en el origen y sus autovalores \(\mu_1\), \(\mu_2\). ¿El origen es un mínimo, un máximo o una silla?
\end{problema}
